\documentclass[12pt, a4paper]{academics}

\academicsCourse{数据库原理与应用}{Database Principle and Application}
\academicsReport{Join 操作查询计划执行}{2026 年 6 月 15 日}
\academicsArthur{华博文}{19240212}

\begin{document}

\def\PYGZus{\textunderscore}
\AtBeginEnvironment{MintedVerbatim}{\catcode`\_=12\relax}
\setminted{mathescape=false}

\makeacademicscover

\section{实验环境}

\subsection{操作系统}

macOS 26.4 arm64，Darwin 25.4.0，Apple M5。

\subsection{数据库平台}

PostgreSQL 16（Docker 容器，\texttt{postgres:16} 镜像）。

\subsection{文档工具}

\LaTeX{}。

\section{实验目的}

\begin{enumerate}
    \item 学习并掌握数据库管理系统中查询处理与查询优化的基本知识和方法——理解查询计划中各算子的含义，以及优化器如何根据统计信息选择连接算法；
    \item 通过 PostgreSQL 的 \texttt{EXPLAIN ANALYZE} 工具，对比 Nested Loop Join、Hash Join、Merge Join 三种连接算法在不同条件下的执行计划与性能表现；
    \item 研究表规模、索引有无等因素对连接算法选择的影响，总结针对不同业务场景提升连接查询性能的策略与方法。
\end{enumerate}

\section{实验内容}

\subsection{环境准备}

启动 Docker 容器并连接：

\begin{minted}{bash}
docker exec -it pg-ex5 psql -U postgres
\end{minted}

\begin{minted}{sql}
\set AUTOCOMMIT off
\end{minted}

\textbf{步骤 1：创建测试数据库与表。}

本实验需要两张等值连接表——小表 \texttt{departments}（10 行）和大表 \texttt{employees}（10\,000 行），通过 \texttt{dept\_id} 关联。后续实验还会构建一张与 employees 等规模的 \texttt{departments\_big} 表。

\begin{minted}{sql}
CREATE DATABASE join_test;
\c join_test

CREATE TABLE departments (
    dept_id   INT PRIMARY KEY,
    dept_name VARCHAR(50)
);

CREATE TABLE employees (
    emp_id    INT PRIMARY KEY,
    emp_name  VARCHAR(50),
    dept_id   INT,
    salary    NUMERIC(10, 2)
);

INSERT INTO departments
SELECT i, 'Department_' || i
FROM generate_series(1, 10) AS i;

INSERT INTO employees
SELECT i,
       'Employee_' || i,
       (random() * 9)::INT + 1,
       (random() * 90000 + 10000)::NUMERIC(10, 2)
FROM generate_series(1, 10000) AS i;

COMMIT;
\end{minted}

\begin{minted}{text}
CREATE DATABASE
You are now connected to database "join_test" as user "postgres".
CREATE TABLE
CREATE TABLE
INSERT 0 10
INSERT 0 10000
COMMIT
\end{minted}

\textbf{步骤 2：确认数据规模。}

\begin{minted}{sql}
SELECT COUNT(*) FROM departments;
SELECT COUNT(*) FROM employees;
\end{minted}

\begin{minted}{text}
 count
-------
    10
(1 row)

 count
-------
 10000
(1 row)
\end{minted}

\begin{figure}[htbp]
\centering
\includegraphics[width=\textwidth]{../res/01-env-setup.png}
\caption{环境准备：建库与建表}
\end{figure}

\subsection{实验一：Nested Loop Join}

Nested Loop Join 以一张表为外表，对每一行扫描内表进行匹配。关键前提是\textbf{内表}连接键上有索引——外表仅需一次 Seq Scan，内表走 Index Scan，避免笛卡尔积。

以下通过 \texttt{SET enable\_hashjoin = off} 和 \texttt{SET enable\_mergejoin = off} 强制优化器选择 Nested Loop Join。

\textbf{步骤 1：仅依赖 departments 主键索引。}employees 上无额外索引，departments 作为内表利用其主键索引。

\begin{minted}{sql}
SET enable_hashjoin  = off;
SET enable_mergejoin = off;

EXPLAIN ANALYZE
SELECT e.emp_name, d.dept_name
FROM employees e
JOIN departments d ON e.dept_id = d.dept_id;
\end{minted}

\begin{minted}{text}
 Nested Loop  (cost=0.16..435.61 rows=10000 width=131)
               (actual time=0.071..7.577 rows=10000 loops=1)
   ->  Seq Scan on employees e  (cost=0.00..184.00 rows=10000 width=17)
                                (actual time=0.012..1.496 rows=10000 loops=1)
   ->  Memoize  (cost=0.16..0.18 rows=1 width=122)
                (actual time=0.000..0.000 rows=1 loops=10000)
         Cache Key: e.dept_id
         Cache Mode: logical
         Hits: 9990  Misses: 10  Evictions: 0  Overflows: 0
         Memory Usage: 2kB
         ->  Index Scan using departments_pkey on departments d
               (cost=0.15..0.17 rows=1 width=122)
               (actual time=0.003..0.003 rows=1 loops=10)
               Index Cond: (dept_id = e.dept_id)
 Planning Time: 0.956 ms
 Execution Time: 8.253 ms
\end{minted}

\textbf{步骤 2：在 employees 上建索引后再次观察。}

\begin{minted}{sql}
CREATE INDEX idx_emp_dept ON employees(dept_id);

EXPLAIN ANALYZE
SELECT e.emp_name, d.dept_name
FROM employees e
JOIN departments d ON e.dept_id = d.dept_id;
\end{minted}

\begin{minted}{text}
 Nested Loop  (cost=0.16..435.61 rows=10000 width=131)
               (actual time=0.023..8.243 rows=10000 loops=1)
   ->  Seq Scan on employees e  (cost=0.00..184.00 rows=10000 width=17)
                                (actual time=0.006..1.443 rows=10000 loops=1)
   ->  Memoize  (cost=0.16..0.18 rows=1 width=122)
                (actual time=0.000..0.000 rows=1 loops=10000)
         Cache Key: e.dept_id
         Cache Mode: logical
         Hits: 9990  Misses: 10  Evictions: 0  Overflows: 0
         Memory Usage: 2kB
         ->  Index Scan using departments_pkey on departments d
               (cost=0.15..0.17 rows=1 width=122)
               (actual time=0.002..0.002 rows=1 loops=10)
               Index Cond: (dept_id = e.dept_id)
 Planning Time: 0.343 ms
 Execution Time: 8.971 ms
\end{minted}

\textbf{结论：}两个步骤的执行计划几乎一致——\texttt{idx\_emp\_dept} 建在 employees（外表）上对 Nested Loop 无影响，因为外表走的是 Seq Scan，Nested Loop 依赖的是\textbf{内表}索引。优化器选择小表 departments（10 行）作为内表，利用其主键索引进行 Index Scan。此外，PostgreSQL 16 的 \texttt{Memoize} 算子缓存了内表查询结果——10\,000 次循环中仅 10 次实际访问索引（misses），其余 9\,990 次全部命中缓存，将实际索引探测开销降至几乎为零。

\begin{figure}[htbp]
\centering
\includegraphics[width=\textwidth]{../res/01-nested-loop.png}
\caption{Nested Loop Join（无额外索引）}
\end{figure}

\begin{figure}[htbp]
\centering
\includegraphics[width=\textwidth]{../res/01-nested-loop-step2.png}
\caption{Nested Loop Join（建 idx\_emp\_dept 后）}
\end{figure}

\subsection{实验二：Hash Join}

Hash Join 先对内表在连接键上构建哈希表，再扫描外表逐行探测。时间复杂度 $O(M + N)$，不依赖索引。

\textbf{步骤 1：小内表（departments，10 行）。}

\begin{minted}{sql}
SET enable_hashjoin  = on;
SET enable_mergejoin = off;
SET enable_nestloop  = off;

EXPLAIN ANALYZE
SELECT e.emp_name, d.dept_name
FROM employees e
JOIN departments d ON e.dept_id = d.dept_id;
\end{minted}

\begin{minted}{text}
 Hash Join  (cost=22.15..232.61 rows=10000 width=131)
             (actual time=0.046..3.887 rows=10000 loops=1)
   Hash Cond: (e.dept_id = d.dept_id)
   ->  Seq Scan on employees e  (cost=0.00..184.00 rows=10000 width=17)
                                (actual time=0.016..1.223 rows=10000 loops=1)
   ->  Hash  (cost=15.40..15.40 rows=540 width=122)
             (actual time=0.018..0.020 rows=10 loops=1)
         Buckets: 1024  Batches: 1  Memory Usage: 9kB
         ->  Seq Scan on departments d  (cost=0.00..15.40 rows=540 width=122)
                                       (actual time=0.009..0.012 rows=10 loops=1)
 Planning Time: 0.269 ms
 Execution Time: 4.434 ms
\end{minted}

\textbf{步骤 2：大内表（departments\_big，10\,000 行）。}

\begin{minted}{sql}
CREATE TABLE departments_big AS
SELECT i AS dept_id, 'Dept_' || i AS dept_name
FROM generate_series(1, 10000) AS i;

EXPLAIN ANALYZE
SELECT e.emp_name, d.dept_name
FROM employees e
JOIN departments_big d ON e.dept_id = d.dept_id;
\end{minted}

\begin{minted}{text}
 Hash Join  (cost=246.88..14682.38 rows=406400 width=45)
             (actual time=3.079..7.973 rows=10000 loops=1)
   Hash Cond: (e.dept_id = d.dept_id)
   ->  Seq Scan on employees e  (cost=0.00..184.00 rows=10000 width=17)
                                (actual time=0.009..1.420 rows=10000 loops=1)
   ->  Hash  (cost=145.28..145.28 rows=8128 width=36)
             (actual time=3.053..3.055 rows=10000 loops=1)
         Buckets: 16384 (originally 8192)  Batches: 1 (originally 1)
         Memory Usage: 597kB
         ->  Seq Scan on departments_big d
               (cost=0.00..145.28 rows=8128 width=36)
               (actual time=0.011..1.444 rows=10000 loops=1)
 Planning Time: 0.196 ms
 Execution Time: 8.644 ms
\end{minted}

\textbf{结论：}Hash Join 在两个规模下均表现良好。小内表（10 行）时哈希表仅 9\,kB，执行时间 4.434\,ms；大内表（10\,000 行）时哈希表 597\,kB，执行时间 8.644\,ms——大部分时间消耗在内表的 Seq Scan 和哈希表构建上，而哈希探测本身极快。对比 Nested Loop（8.253\,ms），小表场景下 Hash Join 反而更快，因为 Hash Join 只扫描两张表各一次，而 Nested Loop 循环 10\,000 次（即便有 Memoize 缓存）。

\begin{figure}[htbp]
\centering
\includegraphics[width=\textwidth]{../res/02-hash-join.png}
\caption{Hash Join（小内表 vs 大内表）}
\end{figure}

\subsection{实验三：Merge Join}

Merge Join 要求两张表在连接键上已排序。若连接键有索引，可直接利用索引顺序，省去显式 Sort；无索引则需先排序再归并。

\textbf{步骤 1：两张表连接键均有索引。}

\begin{minted}{sql}
CREATE INDEX idx_dept_big ON departments_big(dept_id);

SET enable_nestloop  = off;
SET enable_hashjoin  = off;
SET enable_mergejoin = on;

EXPLAIN ANALYZE
SELECT e.emp_name, d.dept_name
FROM employees e
JOIN departments_big d ON e.dept_id = d.dept_id;
\end{minted}

\begin{minted}{text}
 Merge Join  (cost=0.57..8578.85 rows=500000 width=45)
              (actual time=0.086..8.616 rows=10000 loops=1)
   Merge Cond: (e.dept_id = d.dept_id)
   ->  Index Scan using idx_emp_dept on employees e
         (cost=0.29..527.57 rows=10000 width=17)
         (actual time=0.038..3.978 rows=10000 loops=1)
   ->  Materialize  (cost=0.29..551.28 rows=10000 width=36)
                    (actual time=0.039..0.891 rows=10001 loops=1)
         ->  Index Scan using idx_dept_big on departments_big d
               (cost=0.29..526.28 rows=10000 width=36)
               (actual time=0.036..0.039 rows=11 loops=1)
 Planning Time: 0.445 ms
 Execution Time: 9.276 ms
\end{minted}

\begin{figure}[htbp]
\centering
\includegraphics[width=\textwidth]{../res/03-merge-join-step1.png}
\caption{Merge Join（有索引）}
\end{figure}

\textbf{步骤 2：删除索引后再次观察。}

\begin{minted}{sql}
DROP INDEX idx_emp_dept;
DROP INDEX idx_dept_big;

EXPLAIN ANALYZE
SELECT e.emp_name, d.dept_name
FROM employees e
JOIN departments_big d ON e.dept_id = d.dept_id;
\end{minted}

\begin{minted}{text}
 Merge Join  (cost=1676.77..9226.77 rows=500000 width=45)
              (actual time=9.667..14.828 rows=10000 loops=1)
   Merge Cond: (e.dept_id = d.dept_id)
   ->  Sort  (cost=848.39..873.39 rows=10000 width=17)
             (actual time=5.697..7.033 rows=10000 loops=1)
         Sort Key: e.dept_id
         Sort Method: quicksort  Memory: 852kB
         ->  Seq Scan on employees e  (cost=0.00..184.00 rows=10000 width=17)
                                     (actual time=0.016..2.475 rows=10000 loops=1)
   ->  Sort  (cost=828.39..853.39 rows=10000 width=36)
             (actual time=3.960..4.692 rows=10001 loops=1)
         Sort Key: d.dept_id
         Sort Method: quicksort  Memory: 775kB
         ->  Seq Scan on departments_big d
               (cost=0.00..164.00 rows=10000 width=36)
               (actual time=0.010..1.914 rows=10000 loops=1)
 Planning Time: 0.313 ms
 Execution Time: 15.575 ms
\end{minted}

\textbf{结论：}有索引时 Merge Join 直接走 Index Scan（利用索引已排序），无需 Sort，执行时间 9.276\,ms；删除索引后，两张表均需显式 Sort（quicksort，分别消耗 852\,kB 和 775\,kB 内存），执行时间增至 15.575\,ms——排序成本约占总时间的 60\%。这清楚表明 Merge Join 的性能关键在于「数据已排序」这一前提。

\begin{figure}[htbp]
\centering
\includegraphics[width=\textwidth]{../res/03-merge-join.png}
\caption{Merge Join（无索引）}
\end{figure}

\section{三种连接算法对比}

\begin{table}[htbp]
\centering
\small
\caption{三种连接算法对比}
\vspace{8pt}
\resizebox{\textwidth}{!}{%
\begin{tabular}{|c|c|c|c|}
\hline
 & \textbf{Nested Loop} & \textbf{Hash Join} & \textbf{Merge Join} \\
\hline
耗时 & 8.3\,ms & 4.4--8.6\,ms & 9.3--15.6\,ms \\
\hline
复杂度 & $O(MN)$ / $O(M)^{*}$ & $O(M+N)$ & $O(M \log M)^{**}$ / $O(M+N)$ \\
\hline
前提 & 内表索引 & 内表装入内存 & 连接键有序 \\
\hline
关键发现 & 缓存命中 99.9\%\newline 索引须建内表侧 & 小表场景最优\newline 不依赖索引 & 无索引时 Sort\newline 占 60\% 耗时 \\
\hline
\end{tabular}%
}

\vspace{4pt}
{\footnotesize * 有 Memoize 缓存时 \qquad ** 无索引、需显式排序时}
\end{table}

\section{实验关键点总结}

\begin{enumerate}
    \item \textbf{Nested Loop 的内外表选择与索引方向}：PostgreSQL 自动选择小表为内表。索引必须建在内表连接键上才有效——本实验中在 employees（外表）上建 \texttt{idx\_emp\_dept} 对 Nested Loop 完全无效。此外，PostgreSQL 14+ 引入的 Memoize 算子显著降低了重复索引探测的代价。
    \item \textbf{Hash Join 在小表场景的意外优势}：Hash Join 不依赖索引，仅需扫描两张表各一次加哈希探测。本实验中 10 行内表的 Hash Join（4.434\,ms）反而比带 Memoize 的 Nested Loop（8.253\,ms）更快，因为 Nested Loop 仍有 10\,000 次循环开销。
    \item \textbf{Merge Join 的排序代价}：有索引时 Merge Join 走 Index Scan（利用索引排序），无索引时需额外 Sort 且占执行时间的大部分。Merge Join 适合主键-外键等值连接——两条索引同时提供了排序和快速定位。
\end{enumerate}

\section{实验过程归纳总结}

本次实验通过 \texttt{EXPLAIN ANALYZE} 对比了三种连接算法的执行计划。回顾这一过程，有以下三点收获。

\textbf{第一，执行计划是查询优化的核心工具。}\texttt{EXPLAIN ANALYZE} 不仅展示优化器选择的算子组合和代价估计（\texttt{cost}），还给出实际执行时间与行数——两者对比可发现统计信息是否过时。例如 Hash Join 步骤 2 中优化器估计行数为 406\,400，实际仅 10\,000 行，说明 \texttt{departments\_big} 刚由 \texttt{CREATE TABLE AS} 生成后未执行 \texttt{ANALYZE}，统计信息缺失。

\textbf{第二，索引必须建在「正确的一侧」。}实验一中在 employees 上建索引对 Nested Loop 毫无影响——因为优化器将 departments 选为内表，需要的是内表索引（departments 主键已有），外表的索引根本用不上。这纠正了「建索引就能加速 JOIN」的常见误解——索引的方向比有没有索引更关键。

\textbf{第三，理论模型与工程实现的差异。}教材中 Nested Loop 的时间复杂度为 $O(M \times N)$，但 PostgreSQL 16 的 Memoize 将重复内表查找缓存为哈希探测，实际代价远低于理论值。同样，Hash Join 的「内表须装入内存」在本次实验中完全满足（最大仅 597\,kB），但在超大数据集下可能触发磁盘溢写（multiple batches），性能将急剧下降。这些工程细节是纸上推演无法覆盖的。

\section{结论}

本实验在 PostgreSQL 16 平台上通过构造不同规模的数据表、在连接键上设置与删除索引、控制优化器参数（\texttt{SET enable\_*}），分别构造了 Nested Loop Join、Hash Join、Merge Join 三种连接方案，并使用 \texttt{EXPLAIN ANALYZE} 对比了执行计划与性能。实验表明：（1）Nested Loop Join 依赖内表索引，外表索引无效；PostgreSQL 16 的 Memoize 大幅降低重复探测开销；（2）Hash Join 不依赖索引，在 10\,000 行规模下综合性能最优（4--9\,ms）；（3）Merge Join 有索引时可省去排序步骤，无索引时排序成本占主导。在实际开发中，小表 JOIN 可默认信任优化器；大表 JOIN 应根据连接键建对索引——Nested Loop 建在内表侧，Merge Join 建在两侧。

\end{document}
